\documentclass[12pt]{article}
\usepackage[T1]{fontenc}
\usepackage[utf8]{inputenc}
\usepackage{lmodern}
\usepackage{textcomp}
\usepackage{enumitem}
\usepackage{geometry}
\usepackage{graphicx}
\usepackage{amsmath, amssymb}
\geometry{margin=1in}
\begin{document}
\begin{center}
Physics - Undergraduate Level Assessment 16 \
Total Marks: 40 \
Answer all questions.
\end{center}

\section*{Multiple Choice (10 marks)}
\begin{enumerate}[label=\arabic*.]
\item The rest mass of a particle with total energy 5.0 GeV and momentum 4.9 GeV/c is approximately
\begin{enumerate}[label=(\alph*)]
    \item 0.1 GeV/$c^2$
    \item 0.2 GeV/$c^2$
    \item 0.5 GeV/$c^2$
    \item 1.0 GeV/$c^2$
\end{enumerate}
\item The mean kinetic energy of the conduction electrons in metals is ordinarily much higher than kT because
\begin{enumerate}[label=(\alph*)]
    \item electrons have many more degrees of freedom than atoms do
    \item the electrons and the lattice are not in thermal equilibrium
    \item the electrons form a degenerate Fermi gas
    \item electrons in metals are highly relativistic
\end{enumerate}
\item The speed of light inside of a nonmagnetic dielectric material with a dielectric constant of 4.0 is
\begin{enumerate}[label=(\alph*)]
    \item 1.2 x $10^9$ m/s
    \item 3.0 x $10^8$ m/s
    \item 1.5 x $10^8$ m/s
    \item 1.0 x $10^8$ m/s
\end{enumerate}
\item A refracting telescope consists of two converging lenses separated by 100 cm. The eye- piece lens has a focal length of 20 cm. The angular magnification of the telescope is
\begin{enumerate}[label=(\alph*)]
    \item 4
    \item 5
    \item 6
    \item 20
\end{enumerate}
\item A satellite of mass m orbits a planet of mass M in a circular orbit of radius R. The time required for one revolution is
\begin{enumerate}[label=(\alph*)]
    \item independent of M
    \item proportional to $m^(1$/2)
    \item linear in R
    \item proportional to $R^(3$/2)
\end{enumerate}
\end{enumerate}

\section*{True / False (10 marks)}
\textit{Answer True or False}
\begin{enumerate}[label=\arabic*.]
\setcounter{enumi}{5}
\item True or False: During a reversible thermodynamic process, the temperature of the system must remain constant.
\item True or False: An observer moving at 0.50 c will measure the length of a 1.00 m rod to be 0.80 m due to relativistic effects.
\item True or False: The Sun's energy is primarily produced by the fusion of two hydrogen atoms into one helium atom, resulting in significant energy release.
\item True or False: The electric displacement current through a surface is proportional to the rate of change of the magnetic flux through that surface.
\item True or False: The detector will detect exactly 10 photons when 100 photons are sent, due to its quantum efficiency of 0.1.
\end{enumerate}

\section*{Long Form Response (20 marks)}
\textit{Answer each question with a clear and complete explanation.}
\begin{enumerate}[label=\arabic*.]
\setcounter{enumi}{10}
\item Discuss the implications of the universe's temperature change from 12 K to 3 K on the spatial distribution of galaxies, focusing on their relative distances during these periods.
\item Discuss how the Hall coefficient can be used to determine the sign of charge carriers in a doped semiconductor, including the underlying principles and implications for semiconductor behavior.
\end{enumerate}

\vfill
\noindent\textit{End of Paper}
\end{document}